\section{Scope, implementation changes, and proof audit}
\label{sec:scope}

\subsection{Results established in this note}

\begin{center}
\small
\begin{tabular}{>{\raggedright\arraybackslash}p{0.21\textwidth}>{\raggedright\arraybackslash}p{0.29\textwidth}>{\raggedright\arraybackslash}p{0.39\textwidth}}
\toprule
Model and decoder & Reference risk & Proved model range\\
\midrule
Ordinary SBM; regularized block Bernoulli LRT
& Full-edge true-parameter Bernoulli row LRT
& General positive full-rank compact profiles, fixed $K$,
$np\to\infty$ and $np^2\le C_1$; \cref{restore:sbm-theorem}.\\[4pt]
DC-SBM; regularized degree-metric reconstruction and fitted profile LRT
& Full-edge profile LRT with true neighbor parameters
& General positive full-rank bounded-parameter class, fixed $K$,
$np\to\infty$ and $np^2\le C_1$; \cref{restore:dc-theorem}.\\[4pt]
DC-SBM; the same degree-metric procedure
& Full-edge LRT additionally supplied with the target's true degree
& The preceding sparse class plus finite-$n$ equal block exposures;
the classical balanced planted-partition local oracle is included.\\[4pt]
Earlier ordinary free-profile schedule
& Full-edge true-parameter Bernoulli row LRT
& Separate unregularized results \cref{thm:sbm,prop:sbm-margin};
these also allow probabilities outside bounded $np^2$.\\[4pt]
General DC-SBM; any graph algorithm
& Known-target-degree local LRT
& A uniform constant ratio is impossible; \cref{thm:imposs} supplies two
explicit normalized parameter sequences.\\
\bottomrule
\end{tabular}
\end{center}

The algorithmic positive statements are bounds on expected misclassification, with
failure events included. None is merely a statement conditional on correct
initial groups, small fitted-parameter error, a favorable LOO projector,
or successful numerical k-means.

\subsection{What changed relative to the recent experimental code}

The recent ordinary implementation evaluates the full Bernoulli working
criterion on signed $A\Pi_X$, uses finite growing/replacement trials, and
performs the mandatory final E/M step. The earlier theorem
\cref{thm:sbm} analyzes that mathematical schedule. The primary theorem
\cref{thm:regularized} instead uses the specified radius-count initializer
and block Bernoulli fits. It does not identify these two schedules.
For arbitrary fixed $K$ the encoder is parameterized
by $K$; a script hardcoded to three eigenvectors is just its $K=3$ case.
If the signal matrix is indefinite, largest absolute eigenvalues are
required. Selecting largest algebraic eigenvalues is covered only when
all informative population eigenvalues are positive.

For the earlier unregularized DC-SBM construction,
$A\Pi_X$ and $AH_D$ are different reconstructions.
\Cref{thm:dc} uses the degree-metric action
$AH_D$ and its explicit hollow correction. It must not be retrospectively
attached to an experimental curve computed with the prior Euclidean
action. The retained input $(D,U,\Lambda)$ is identical; the decoder
changes. The radius-count initializer is a concrete certified alternative
to unquantified fixed-restart Lloyd clustering. A certified constant-factor
k-means initializer could also be used, but its own failure probability
would have to be paid at the oracle scale.

The regularized algorithm of \cref{restore:algorithm} recomputes the
eigensystem with $G_i=\max(D_i,a_0\bar D)$, then uses an exact rank-one
degree correction and the simple hollowing in
\eqref{restore:reconstruction}. Its clamped DC nuisance estimates remain
bounded even when individual raw degrees are exceptional.
An isolated vertex affects its own row without forcing a whole-graph
fallback. These are actual algorithm specifications, not modifications
made only to a population surrogate. Earlier experimental curves do
not automatically evaluate this regularized variant.

The theorems use exact algebraic spaces and exact arithmetic, including
the specified exact E/M updates for the ordinary decoder and the stated
block fits and target-profile maxima for DC-SBM. The ordinary finite
schedule does not assume a global maximizer of the mixture likelihood.
They do not assert that a fixed floating-point tolerance remains adequate
as $n\to\infty$. In particular, augmenting $D^{-1/2}U$ by an unscaled
constant column and applying a fixed relative rank cutoff is not an
asymptotic definition of its algebraic rank. Because the constant vector is
already retained, it can be omitted from that numerical augmentation, or
columns can be scaled before orthogonalization. Quantified numerical
errors may be admitted if their induced likelihood-contrast and
cumulant errors are uniformly bounded.

\subsection{Why this is a risk bound}

The error incurred in replacing a full-graph spectral score by its LOO
version is deterministically $O(1)$ on the stated aggregate and local
target events. Fitted block parameters cost $O(1)$ in both new
decoders, uniformly over the invariant labeling basin.
In the separate earlier ordinary proof, free profiles produce
$O(1+\mathcal B_i)$ with a bounded exponential moment.
The remaining projection error against the raw oracle row is random:
it is not asserted to be uniformly bounded for every possible row.
Its mean under the oracle error-boundary tilt and its tilted quadratic
moment are bounded. Exceptional large coefficients are paid only when
the independent tested row encounters them.
\Cref{restore:tail-energy-lemma,lem:sparse-truncated-transfer}
convert these facts into a constant risk multiplier while preserving
$d^{-1/2}e^{-I}$. No simultaneous row-incoherence event is required
by the regularized theorem.

The weighted cancellation controls the first-order mean term that a
generic unweighted dot product need not control in DC-SBM. It does not
assert exact recovery of every individual edge or of every observed
block count. The Bernoulli profile reduction has a visible $O(np^2)$
remainder; extending the DC theorem beyond bounded $np^2$ requires a new
argument for that remainder and the heterogeneous nonedge weights.

The regularized algorithmic theorem covers $np=o(\log n)$ whenever
$np\to\infty$, as well as every higher degree scale with bounded
$np^2$. A vanishing normalized signal singular value and unbounded
$np^2$ in the general DC theorem remain outside its stated class.
An unrestricted known-$\theta_i$ oracle in general DC-SBM is
explicitly impossible by \cref{thm:imposs}; this obstruction does not
apply to the equal-exposure result. The unchanged ordinary
free-profile schedule retains the scope of its separate theorem.
\Cref{sec:sparse-extension} supplies the regularized geometry and
sharp-transfer tools, and \cref{sec:sparse-restored} completes their
row-tail, local LOO, fitting, and unconditional probability interfaces.

\subsection{Audit record}

The generalized-eigenvector full-to-LOO derivation was checked independently
against a second argument based on a normalized arrowhead decomposition
and a residual estimate for the degree-renormalized minor.
The row-incoherence bootstrap uses a verified weighted Bernoulli
concentration statement and controls every deleted basis simultaneously.
The finite ordinary likelihood schedule was checked against the current
implementation, including signed observations, trial resets, strict
replacement acceptance and the mandatory endpoint order.
The DC proof was separately checked for the uniform effect of
data-dependent label errors, exact Bernoulli block fitting, the target
profile cap, degree-zero oracle rows and adaptive finite hard refits.
The two-point obstruction was verified independently from its affinity
and permutation-loss calculations.

For the sparse restoration, independent checks covered the
same-threshold row-energy recurrence with square-root degree caps,
the empirical-threshold perturbation, local generalized eigenvector
equations, weighted wrong-label mass, signed Bernoulli fitting,
finite hard-refit invariance, and the ordinary unprofiled intercept.
The high-scale probability closure was checked directly against
the centered-entry norm inequality of Bandeira--van Handel.
Training events are measurable from their deleted graph; full-graph
events are used only for deterministic error inclusions.

These are internal mathematical and implementation checks, not external
peer review. The note provides the complete proofs so each claim can be
reviewed directly. There is no simulation estimate being used as proof
of a bounded asymptotic ratio.

